\chapter{Agents LLM et utilisation d'outils}
\label{ch:agents}

\begin{quote}
\emph{<< Un agent est un LLM qui peut prendre des actions, pas seulement g\'en\'erer du texte. Il raisonne sur ce qu'il faut faire, utilise des outils, et it\`ere jusqu'\`a r\'esoudre la t\^ache. >>}
\end{quote}

% --- hook ---
Un LLM standard prend un prompt et renvoie du texte. Un agent prend un prompt et \emph{agit} --- il interroge une API, lit un fichier, exécute du code, observe le résultat, et recommence jusqu'à avoir résolu la tâche. Même modèle de base, mais avec un échafaudage qui transforme un générateur de texte en système autonome. C'est aussi là que les enjeux de sûreté deviennent concrets : un agent qui se trompe ne dit pas une bêtise, il l'exécute.

% ============================================================
\section{Qu'est-ce qu'un agent LLM ?}

Un LLM standard prend un prompt et renvoie du texte. Un \textbf{agent} ajoute :
\begin{itemize}
  \item \textbf{Outils} : fonctions que le LLM peut appeler (recherche, calculatrice, ex\'ecution de code, API).
  \item \textbf{Raisonnement} : le LLM d\'ecide quel outil utiliser et dans quel ordre.
  \item \textbf{M\'emoire} : l'agent maintient le contexte \`a travers plusieurs \'etapes.
  \item \textbf{It\'eration} : l'agent boucle jusqu'\`a r\'esoudre la t\^ache ou atteindre une condition d'arr\^et.
\end{itemize}

% ============================================================
\section{Le pattern ReAct}

ReAct (Reasoning + Acting) entrelace les traces de raisonnement avec les actions sur outils :

\begin{enumerate}
  \item \textbf{Thought} : le LLM raisonne sur l'\'etat courant et la prochaine \'etape.
  \item \textbf{Action} : le LLM appelle un outil avec des entr\'ees sp\'ecifiques.
  \item \textbf{Observation} : l'outil renvoie un r\'esultat.
  \item R\'ep\'eter jusqu'\`a ce que le LLM ait assez d'information pour r\'epondre.
\end{enumerate}

\begin{minted}{python}
# Prompt de style ReAct (conceptuel)
react_prompt = """Answer the question using the available tools.

Tools:
- search(query): search the web for information
- calculator(expression): evaluate a math expression

Question: What is the population of Benin multiplied by 3?

Thought: I need to find the population of Benin first.
Action: search("population of Benin 2025")
Observation: The population of Benin is approximately 14.4 million.
Thought: Now I need to multiply 14.4 million by 3.
Action: calculator("14400000 * 3")
Observation: 43200000
Thought: I have the answer.
Answer: The population of Benin multiplied by 3 is approximately 43,200,000."""

print(react_prompt)
\end{minted}

% ============================================================
\section{Function calling avec Groq}

Les API LLM modernes prennent en charge des appels d'outils/fonctions structur\'es :

\begin{minted}{python}
from groq import Groq
import json

client = Groq()

# D\'efinir les outils
tools = [
    {
        "type": "function",
        "function": {
            "name": "get_weather",
            "description": "Get the current weather for a city",
            "parameters": {
                "type": "object",
                "properties": {
                    "city": {"type": "string", "description": "City name"},
                    "unit": {"type": "string", "enum": ["celsius", "fahrenheit"]},
                },
                "required": ["city"],
            },
        },
    },
    {
        "type": "function",
        "function": {
            "name": "calculate",
            "description": "Evaluate a mathematical expression",
            "parameters": {
                "type": "object",
                "properties": {
                    "expression": {"type": "string", "description": "Math expression"},
                },
                "required": ["expression"],
            },
        },
    },
]

# Impl\'ementations effectives des outils
def get_weather(city, unit="celsius"):
    # En production, appeler une vraie API m\'et\'eo
    return {"city": city, "temperature": 28, "unit": unit, "condition": "sunny"}

def calculate(expression):
    return {"result": eval(expression)}  # attention : eval est dangereux en production

tool_map = {"get_weather": get_weather, "calculate": calculate}

# Boucle de l'agent
messages = [{"role": "user", "content": "What is the temperature in Cotonou in Fahrenheit?"}]

response = client.chat.completions.create(
    model="llama-3.1-8b-instant",
    messages=messages,
    tools=tools,
    tool_choice="auto",
)

# Traiter les appels d'outils
msg = response.choices[0].message
if msg.tool_calls:
    messages.append(msg)
    for tc in msg.tool_calls:
        func_name = tc.function.name
        args = json.loads(tc.function.arguments)
        result = tool_map[func_name](**args)
        messages.append({
            "role": "tool",
            "tool_call_id": tc.id,
            "content": json.dumps(result),
        })
    # Obtenir la r\'eponse finale
    final = client.chat.completions.create(
        model="llama-3.1-8b-instant", messages=messages
    )
    print(final.choices[0].message.content)
\end{minted}

% ============================================================
\section{Agents LangChain}

LangChain fournit un cadre d'agent de plus haut niveau :

\begin{minted}{python}
from langchain_google_genai import ChatGoogleGenerativeAI
from langchain_core.tools import tool
from langchain.agents import create_tool_calling_agent, AgentExecutor
from langchain_core.prompts import ChatPromptTemplate

# D\'efinir les outils
@tool
def search_arxiv(query: str) -> str:
    """Search arxiv.org for recent papers. Returns titles and abstracts."""
    import urllib.request, json
    url = f"http://export.arxiv.org/api/query?search_query=all:{query}&max_results=3"
    response = urllib.request.urlopen(url).read().decode()
    # Parsing simplifi\'e
    return response[:2000]

@tool
def python_calculator(expression: str) -> str:
    """Evaluate a Python math expression. Example: '2**10 + 3*4'"""
    try:
        result = eval(expression)
        return str(result)
    except Exception as e:
        return f"Error: {e}"

# Cr\'eer l'agent
llm = ChatGoogleGenerativeAI(model="gemini-2.0-flash", temperature=0)
tools_list = [search_arxiv, python_calculator]

prompt = ChatPromptTemplate.from_messages([
    ("system", "You are a helpful research assistant. Use tools when needed."),
    ("human", "{input}"),
    ("placeholder", "{agent_scratchpad}"),
])

agent = create_tool_calling_agent(llm, tools_list, prompt)
executor = AgentExecutor(agent=agent, tools=tools_list, verbose=True)

result = executor.invoke({"input": "Find recent papers about LoRA fine-tuning and tell me how many authors the first paper has."})
print(result["output"])
\end{minted}

% ============================================================
\section{Agents \`a raisonnement multi-\'etapes}

Les agents peuvent encha\^iner plusieurs appels d'outils pour r\'esoudre des t\^aches complexes :

\begin{minted}{python}
@tool
def read_file(path: str) -> str:
    """Read the contents of a text file."""
    with open(path) as f:
        return f.read()[:5000]

@tool
def summarize_text(text: str) -> str:
    """Summarize a long text into 3 bullet points."""
    llm = ChatGoogleGenerativeAI(model="gemini-2.0-flash", temperature=0.3)
    response = llm.invoke(f"Summarize in 3 bullet points:\n{text}")
    return response.content

@tool
def write_file(path: str, content: str) -> str:
    """Write content to a file."""
    with open(path, "w") as f:
        f.write(content)
    return f"Written to {path}"

tools_list = [read_file, summarize_text, write_file]
agent = create_tool_calling_agent(llm, tools_list, prompt)
executor = AgentExecutor(agent=agent, tools=tools_list, verbose=True)

result = executor.invoke({
    "input": "Read notes.txt, summarize it, and save the summary to summary.txt"
})
\end{minted}

% ============================================================
\section{Bases de LangGraph}

LangGraph permet de construire des flux de travail d'agents stateful et multi-\'etapes comme des graphes :

\begin{minted}{python}
from langgraph.graph import StateGraph, END
from typing import TypedDict, Annotated

class AgentState(TypedDict):
    question: str
    research: str
    draft: str
    review: str
    final_answer: str

def research_step(state: AgentState) -> AgentState:
    """Research the question using the LLM."""
    llm = ChatGoogleGenerativeAI(model="gemini-2.0-flash")
    result = llm.invoke(f"Research this topic thoroughly: {state['question']}")
    return {"research": result.content}

def draft_step(state: AgentState) -> AgentState:
    """Write a draft answer based on research."""
    llm = ChatGoogleGenerativeAI(model="gemini-2.0-flash")
    result = llm.invoke(
        f"Based on this research:\n{state['research']}\n\n"
        f"Write a clear, concise answer to: {state['question']}"
    )
    return {"draft": result.content}

def review_step(state: AgentState) -> AgentState:
    """Review and improve the draft."""
    llm = ChatGoogleGenerativeAI(model="gemini-2.0-flash")
    result = llm.invoke(
        f"Review this draft for accuracy and clarity. Suggest improvements:\n"
        f"{state['draft']}"
    )
    return {"review": result.content}

def finalize_step(state: AgentState) -> AgentState:
    """Produce the final answer incorporating review feedback."""
    llm = ChatGoogleGenerativeAI(model="gemini-2.0-flash")
    result = llm.invoke(
        f"Original draft:\n{state['draft']}\n\n"
        f"Review feedback:\n{state['review']}\n\n"
        f"Write the final polished answer."
    )
    return {"final_answer": result.content}

# Construire le graphe
workflow = StateGraph(AgentState)
workflow.add_node("research", research_step)
workflow.add_node("draft", draft_step)
workflow.add_node("review", review_step)
workflow.add_node("finalize", finalize_step)

workflow.set_entry_point("research")
workflow.add_edge("research", "draft")
workflow.add_edge("draft", "review")
workflow.add_edge("review", "finalize")
workflow.add_edge("finalize", END)

app = workflow.compile()

result = app.invoke({"question": "What are the key differences between LoRA and QLoRA?"})
print(result["final_answer"])
\end{minted}

\begin{aitip}
LangGraph est plus puissant que les agents LangChain simples pour les flux qui demandent du branchement conditionnel, des boucles, ou des points de contr\^ole human-in-the-loop. Utilisez les agents LangChain pour de simples appels d'outils ; utilisez LangGraph pour des workflows multi-\'etapes complexes.
\end{aitip}

\begin{exercise}
\begin{enumerate}
  \item Construisez un agent LangChain avec 3 outils : recherche web, calculatrice, date/heure courante. Testez-le sur 5 requ\^etes qui demandent des combinaisons d'outils diff\'erentes.
  \item Impl\'ementez le function calling avec l'API Groq. Cr\'eez un outil qui interroge des informations sur un pays (capitale, population, PIB) et un agent qui r\'epond \`a des questions de g\'eographie.
  \item Construisez un workflow LangGraph avec 4 \'etapes : (1) recherche, (2) plan, (3) r\'edaction, (4) revue. Utilisez-le pour g\'en\'erer un court article sur n'importe quel sujet.
  \item Cr\'eez un agent << assistant de code >> qui peut \'ecrire du Python, l'ex\'ecuter, et d\'eboguer les erreurs automatiquement.
  \item Comparez le comportement d'un agent LangChain simple vs d'un workflow LangGraph sur la m\^eme t\^ache complexe. Lequel produit de meilleurs r\'esultats ?
\end{enumerate}
\end{exercise}

\begin{ethicsbox}
Les agents peuvent prendre des actions r\'eelles : envoyer des e-mails, modifier des fichiers, ex\'ecuter du code, appeler des API. Cela amplifie autant les b\'en\'efices que les risques de l'IA. Mettez en place ces garde-fous :
\begin{itemize}
  \item \textbf{Moindre privil\`ege :} ne donner aux agents que les outils dont ils ont besoin.
  \item \textbf{Approbation humaine :} exiger une confirmation humaine pour les actions \`a fort enjeu (envoyer un e-mail, supprimer des donn\'ees).
  \item \textbf{Sandboxing :} ex\'ecuter le code dans des environnements isol\'es.
  \item \textbf{Rate limiting :} emp\^echer les agents de faire un nombre illimit\'e d'appels API ou d'actions.
  \item \textbf{Journalisation d'audit :} consigner chaque action de l'agent pour revue.
\end{itemize}
\end{ethicsbox}

% ============================================================
\section{R\'esum\'e du chapitre}

\begin{itemize}
  \item Un agent LLM combine raisonnement, utilisation d'outils, m\'emoire et it\'eration.
  \item Le pattern ReAct entrelace pens\'ee et action dans une boucle.
  \item Le function calling permet aux LLM d'invoquer des outils structur\'es via API.
  \item Les agents LangChain fournissent un cadre de haut niveau pour les LLM utilisateurs d'outils.
  \item LangGraph permet des workflows stateful multi-\'etapes avec branchement conditionnel.
  \item La s\^uret\'e des agents demande moindre privil\`ege, supervision humaine, sandboxing et journalisation d'audit.
\end{itemize}
